\section{问题重述与分析}
热风烘干通过环境与药材之间的热量和水分交换改变材料内部状态。药材内部不同位置的升温与失水速度并不相同，因而仅用烘房温度或整体平均含水率，不能直接描述轴心的干燥进程。题目给出长25\unit{cm}、半径2\unit{cm}的近似圆柱药材，初始温度28\celsius、干基含水率2.55\unit{kg/kg}，并提供环境记录、半径变化记录及分问物性参数\cite{problem}。

四个问题构成逐步增加模型机制的计算任务：第一问采用附录2描述前1800\unit{s}的预热过程，输出指定时刻及径向位置的温度与含水率；第二问统一采用附录3经验式建立变物性模型，其论文表明确取前3\unit{h}内的分布；第三问以药材各处含水率均低于0.15\unit{kg/kg}为结束条件，确定干燥时长；第四问结合附件2及附录4，进一步考虑几何收缩对干燥时长的影响。本稿完成第一问与第二问。第二问逐秒文件也采用前3\unit{h}这一已披露的工作范围；题面“整个烘干过程”的长期预测与逐秒文件终点不能混同，第三、第四问的长期环境延拓、达标判据与移动边界仍需分别验证。

前两问的输入是烘房环境记录，并没有药材内部温湿响应的观测数据，因此采用基于守恒关系的传输模型，而不进行内部响应拟合。取圆柱轴向中点为$z=0$，将“到药材中心的距离”解释为中截面内到轴线的径向距离$r$，输出不解释为端面局部值或全长平均值。长度与直径之比为6.25；至1800\unit{s}，热扩散尺度约1.74\unit{cm}，远小于中截面距端面的12.5\unit{cm}，为径向近似提供尺度依据。已有有限圆柱二维导热对照在七个规定时刻、中截面21个径向位置与一维解的差异约为$10^{-7}\unit{K}$量级，与所作模态截断对照的变化同量级。这一检验支持第一问时段的中截面温度近似，不构成端面或二维非线性水分场的误差保证；第二问另以有限圆柱耦合计算检查延长时段后的几何近似。

\section{第一问的模型假设与符号}
\subsection{模型假设及其适用范围}
\begin{enumerate}[label=（\arabic*）]
\item 将药材视为均质、各向同性的固定圆柱，初态在空间均匀；本问保持$R=0.02\unit{m}$、$L=0.25\unit{m}$，几何收缩留待第四问处理。固定几何属于分问模型约定，不能由预热时间较短推得真实收缩必然很小。
\item 内部采用单一局部温度，热储存项使用附录2给定的固定有效体积热容$\rho c_p$，导热系数不随状态变化。$\rho=820\unit{kg/m^3}$在本问作为热方程参考参数使用，不同时解释为失水期间始终不变的真实湿体密度。
\item 干物质总量保持不变，固定骨架下的干物质密度$\rho_d$均匀且不变；局部含水率按附录2决定$D(C)$。水分守恒方程中约去$\rho_d$，求解干基含水率本身不需要另行指定干密度数值。
\item 附件1给出的环境是表面外侧空间均匀的已知驱动，采用题给常数换热系数和等效传质系数。对时间记录作分段线性插值，不把末时刻环境替换为整个阶段的恒定边界。
\item 表面水分交换采用题意等效尺度下的含水率差Robin条件。药材干基含水率以干物质质量为分母；若空气量按湿度比解释，则以干空气质量为分母。两者单位同为kg/kg，并不意味着定义相同。所用边界是本问的等效闭合关系，不是已经标定的气固平衡关系。
\item 主热模型计入显热传导和表面对流，不显式加入相变潜热、壁面辐射、支撑接触传热及内部热源。这些机制的省略用于构成题给数据可闭合的基础模型，其影响在后文作条件性讨论，不预先判定均可忽略。
\end{enumerate}

\subsection{主要符号与参数}
\begin{table}[htbp]
\centering\small
\caption{第一问主要符号、参数及变量口径}\label{tab:symbols}
\begin{tabularx}{\textwidth}{@{}lXl@{}}
\toprule
符号 & 含义 & 数值或单位\\
\midrule
$r,t$ & 径向坐标、时间 & m，s\\
$R,L$ & 圆柱半径、长度 & 0.02 m，0.25 m\\
$T,T_s,T_\infty$ & 局部、表面、烘房温度 & $^\circ$C\\
$C,C_s$ & 局部、表面干基含水率 & kg/kg\\
$C_\infty$ & 环境量在等效边界关系中的输入数值 & kg/kg\\
$\rho,c_p$ & 参考密度、比热容 & 820 kg/m$^3$；2600 J/(kg$\cdot$K)\\
$k,h_T$ & 导热系数、对流换热系数 & 0.36 W/(m$\cdot$K)；25 W/(m$^2\cdot$K)\\
$h_m$ & 等效对流传质系数 & $8\times10^{-7}$ m/s\\
$D(C)$ & 局部有效水分扩散系数 & m$^2$/s\\
$\rho_d$ & 固定体积内的干物质质量密度 & kg干物质/m$^3$\\
$\alpha$ & 热扩散率，$k/(\rho c_p)$ & m$^2$/s\\
$\overline T,\overline C$ & 中截面按面积加权的平均状态 & $^\circ$C，kg/kg\\
\bottomrule
\end{tabularx}
\end{table}
\FloatBarrier
